\subsection{Conformal Prediction}
\label{sec2.3: conformal prediction}
\noindent Conformal prediction, initially developed by Vovk, Gammerman and Shafer \cite{vovk2005algorithmic} and then later made
accessible by Angelopoulos and Bates \cite{angelopoulos2023conformal}, is a framework for constructing prediction sets that 
have a coverage guarantee. Unlike the methods seen in \ref{sec2.2: current approaches} it can provide us with a set of possible
outcomes (not a point prediction), have a mathematical guarantee and yet requires no distributional assumptions on the data, 
except for the exchangibility condition. In this section we will develop the framework of conformal prediction and see 
\todo{complete this sentence}

\begin{definition}[Exchangeability]
\label{def: exchangeability}
A finite sequence of random variables $Z_1, Z_2, \dots, Z_m$ taking values in a space $\mathcal{Z}$ is exchangeable if, 
for every permutation $\sigma$ of the indices $\{1, 2, \dots, m\}$, the joint distribution of the permuted sequence is 
identical to that of the original sequence:
$$(Z_1, Z_2, \dots, Z_m) \overset{d}{=} (Z_{\sigma(1)}, Z_{\sigma(2)}, \dots, Z_{\sigma(m)})$$
where $\overset{d}{=}$ denotes equality in distribution.

\noindent Or, equivalently:
$$\mathbb{P}(Z_1 \in A_1, Z_2 \in A_2, \dots, Z_m \in A_m) = \mathbb{P}(Z_{\sigma(1)} \in A_1, Z_{\sigma(2)} \in A_2, 
\dots, Z_{\sigma(m)} \in A_m)$$ for all measurable sets $A_1, A_2, \dots, A_m \subseteq \mathcal{Z}$ and all permutations 
$\sigma \in S_m$ (where $S_m$ is the symmetric group on $m$ elements).
\end{definition}

\vspace{0.15cm}
\noindent Independent and identically distributed (i.i.d.) seqeunces are always exchangeable, but the converse does 
not always hold. Thus, exchangeability is considered to be a weaker condition than i.i.d., since it requires that the 
observations have no special ordering. This condition can easily be met by splitting any dataset. \cite{ochoarivera2024conformal, 
vovk2005algorithmic}  

\subsubsection{Mathematical Setup and Prediction Set}
\label{sec2.3.1: setup}
% \noindent - framework\\
% - coverage guarantee\\
% - exchnagebility
% state Theorem 2.1 formally with both bounds ???
\noindent Let $(X,Y) \in \mathcal{X} \times \mathcal{Y}$ be a random input-label pair, where $\mathcal{X}$ is the feature space that 
represents our object of interest (in our case it is the phage protein level score vectors) and $\mathcal{Y} = \{1, \cdots, K\}$
is the discrete label space (set of candidate hosts). 

\vspace{0.15cm}
\noindent We assume that we have a calibration dataset, $D_{\text{cal}} = \{(X_1, Y_1), \cdots , (X_n, Y_n)\}$, of $n$ 
exchangable data points drawn from an unknown joint distribution, $P_{X,Y}$. Then, for a new unlabelled test input, $X_{n+1}$
our aim is to to construct a prediction set function $C: \mathcal{X} \rightarrow 2^{\mathcal{Y}}$ dependening on the data
which maps the test input to a subset of candidate labels, i.e., $C(X_{n+1}) \subseteq \mathcal{Y}$, such that the true label 
falls in in the set with a defined coverage that is:
$$\mathbb{P}(Y_{n+1} \in C(X_{n+1})) \geq 1 - \alpha$$
where $\alpha \in (0,1)$ is the significance level. In our experiments we set $\alpha = 0.1$  so that the guarantee is 
$90\%$.

\vspace{0.15cm}
\noindent A prediction set function can provide us three qualitatively different outcomes for any new test point:
\begin{enumerate}[label=(\roman*)]
    \item $|C(x)| = 0$ (Abstension): this means that the model cannot find any plausibl host. This happens when the score 
           profile of the text phage is very different from anything seen while calibrating.
    \item $|C(x)| = 1$ (Confident Prediction): this means that exactly one host label is consistent with the calibration
          distribution at $1-\alpha$ significance level.
    \item $|C(x)| \geq 2$ (Unsure): this means that according to the model there are several plausible hosts.  
\end{enumerate}

\noindent The expected set size, $\mathbb{E}[|C(x)|]$ can help us determine the efficiency of the model. If the coverage
is maintained, then the smaller the average set size the more efficient the model is, as it provides us with a more informative
prediction. An ideal predictor should acheive a coverage $\geq 1-\alpha$ with the smallest possible set size. This is the goal
that we will work on in \todo{refer to the correct section}.


\subsubsection{Split conformal prediction} 
\label{sec2.3.2: split conformal}
% \noindent - why we use it, the train/calibration/test separation, validity still holds
\noindent In this work, we have used split conformal prediction, in this the training data is split into two sets: 
a training set and a calibration set. 

\vspace{0.15cm}
\noindent The first set, the training set, is used to fit a score model $\hat f: \mathcal{X} \times \mathcal{Y \rightarrow 
\mathbb{R}}$. We then define a nonconformity score function: $s: \mathcal{X} \times \mathcal{Y} \to \mathbb{R}$ based on 
$\hat{f}$. This score function is used to calculate the error margins, it quantifies how unusual the pair $(x, y)$ is 
relative to the training data.  

\vspace{0.15cm}
\noindent Then the calibration set (the second set), $\{(X_i, Y_i)\}^n_{i=1}$, which is seperate from the training set, is 
used to evaluate the nonconformity scores:
$$S_i = S(X_i, Y_i), \quad \forall i = 1, \cdots, n$$

\vspace{0.15cm}
\noindent To guranatee that out error rate is less than $\alpha$ we need to compute the conformal calibration qunatile,
$\hat q$, from these calibration scores: 

\begin{definition} [Conformal Quantile]
The conformal calibration quantile $q_{\alpha}$ is defined as the empirical $\frac{\lceil (n+1)(1-\alpha) \rceil}{n}$-th 
quantile of the sorted calibration scores $\{S_1, \dots, S_n\}$:$$q_{\alpha} := \text{Quantile}\left( \{S_1, \dots, S_n\}, ; 
\frac{\lceil (n+1)(1-\alpha) \rceil}{n} \right)$$
Alternatively, using the emprical distribution function: 
$$q_{\alpha} = \inf \left\{ q \in \mathbb{R} : \frac{1}{n}\sum_{i=1}^{n} \mathbf{1}(S_i \leq q) \geq \frac{\lceil (n+1)
(1-\alpha) \rceil}{n} \right\}.$$
\end{definition}

\vspace{0.15cm}
\noindent Then, for a new input, $X_{n+1}$ the prediction set is constructed by collecting all the candidate labels in
$y \in \mathcal{Y}$ whose nonconformity scores do not exceed this calibrated threshold:
\begin{equation}
\label{eq: prediction set definition}
C(X_{n+1}) = \left\{ y \in \mathcal{Y} \;\Big|\; s(X_{n+1}, y) \le q_{\alpha} \right\}
\end{equation}


\noindent By the exchangibility assumption \ref{def: exchangeability}, we get that the test score, $S_{n+1} = S \left(X_{n+1}, 
Y_{n+1}\right)$ has rank which is uniformly distributed among $\{S_1, \cdots S_n, S_{n+1}\}$. We get the factor of 
$\frac{n+1}{n}$ due to the finite-sample exchangeability assumption of the $n$ calibration points added to the $1$ 
unseen test point, the probability is split across $n+1$ equal segments. Multiplying $(1-\alpha)$ by $\frac{n+1}{n}$ 
guarantees exact finite-sample coverage \cite{lei2018distribution}: 
$$\mathbb{P}\left(Y_{n+1} \in C(X_{n+1})\right) = \mathbb{P}(S_{n+1} \leq q_{\alpha}) \ge 1 - \alpha.$$ 

\begin{theorem}[Marginal Coverage Guarantee \cite{vovk2005algorithmic}]
\label{thm: marginal coverage}
If the data points $Z_1, \dots, Z_{n+1}$ are exchangeable, then the prediction set $C(X_{n+1})$ defined in 
Equation~\eqref{eq: prediction set definition} satisfies:
\begin{equation*}
\mathbb{P}\left( Y_{n+1} \in C(X_{n+1}) \right) \ge 1 - \alpha
\end{equation*}
Furthermore, if the nonconformity scores $V_i$ are almost surely distinct (continuous random variables), the coverage 
bound is tight:
\begin{equation*}
1 - \alpha \le \mathbb{P}\left( Y_{n+1} \in C(X_{n+1}) \right) \le 1 - \alpha + \frac{1}{n+1}
\end{equation*}
\end{theorem}


\vspace{0.15cm}
\noindent This mathematical guarantee in Theorem \ref{thm: marginal coverage} holds exactly in finite samples without 
any assumptions on the nonconformity score function, $s$, or on the distribution of $(X, Y)$. The choice of $S$ affects 
the efficiency as it impacts how tight the prediction sets are, but it does not effect the validity of the method. 
This makes seperation makes conformal prediction powerfuls, as it enables us to use any score function and yet have the 
coverage guarantee \cite{ochoarivera2024conformal}.

\noindent Overall, this is a valid technique that achieves the coverage guarantee over the calibration set in a single pass 
which makes it computationally efficient.


\subsubsection{Marginal \& Conditional Coverage}
\label{sec2.3.3: coverage types}
\noindent Standard conformal prediction satisfies marginal coverage, which means that the $(1-\alpha)\%$ coverage guarantee holds
on average across all possible choices of the calibration set, $D_{\text{cal}}$, and test samples, $(X_{n+1}, Y_{n+1})$:
$$\mathbb{P}(Y_{n+1} \in C(X_{n+1})) \geq 1-\alpha.$$

\noindent We can also express this as an expectation over both $D_{\text{cal}}$ and the test pair as:
$$\mathbb{E}_{D_{cal}, (X_{n+1}, Y_{n+1})} \left[ \mathbb{I}(Y_{n+1} \in C(X_{n+1})) \right] \geq 1 - \alpha$$

\vspace{0.15cm}
\noindent Even though the marginal coverage is a powerful guarantee but it does not ensure the conditional coverage,
$\mathbb{P}(Y_{n+1}\in C(X_{n+1}) \mid X_{n+1} = x) \ge 1 - \alpha$ for all $x$. Also, for all different classes, $y$
it cannot ensure class conditional coverage: $\mathbb{P}(Y_{n+1} \in C(X_{n+1}) \mid Y_{n+1} = y) \geq 1 - \alpha$. So, for
a dataset that has class imbalance, which is very common in phage-host interaction, a model that only gives is a marginal 
coverage guarantee can potentially under-cover a rare small class while over cover the majority class and still maintain 
the average coverage \cite{ochoarivera2024conformal, romano2020classification}.. 

\vspace{0.15cm}
\noindent For example, suppose we have $\mathcal{y} = \{-1, 1\}$, with $\mathbb{P}(Y = 1) = 0.1$ and $\mathbb{P}(Y=-1) = 0.9.$
Then a predictor that acheives $100\%$ coverage for $Y = -1$ and $0\%$ for $Y = 1$ will still manage to satisfy the 
marginal coverage of $90\%$. So,  in a phage=hot classification setting , where there is asymmtery in the calibration set
sizes for different hosts, a predictor may undercover the minority class below the $1-\alpha$ target coverage. 

\vspace{0.15cm}
\noindent This limitation motivates the class-conditional calibration coverage structure in \todo{refer to section}. In
this framework, originally introduced under Mondrian conformal prediction \cite{vovk2005algorithmic}, the guarantee hold 
for each class individually:
$$\mathbb{P}(Y \in C(X) \mid Y = y) \geq 1-\alpha \quad \forall y \in \mathcal{y}$$

\noindent It is impossible to get the feature-conditional coverage ($X = x$) without additional assumptions, as it would require the
conditional coverage at every point $x$, which is impossible with finite data \cite{ochoarivera2024conformal}. However, we
can get the class-conditional coverage ($Y = y$) is exact and valid for finite-sample. We can compute a seperate 
threshold for each class using only the calibration points with the true label $y$. This is like class conditional 
calibration, which is a justified approximation \todo{justify it mathematically}. 

\vspace{0.15cm}
\noindent This technique has been used in this report and is formalised in \todo{refer to section}. The main point here is
that this class conditional calibration still uses the exact same quantile argument as the marginal calibration, it just 
applies it seperately to each class. Therefore, it is mathematically valid under the assumption of exchangeability
within each class.
